\exercisesection{}

在本节习题中，除非另有说明，G 表示赋有左哈尔测度的局部紧拓扑群，该测度记为 \(\mu\)。

\begin{enumerate}
\def\labelenumi{\arabic{enumi})}
\item
  令 \(\pi\) 为 G 在复希尔伯特空间 E 上的酉表示。表示 \(\pi\) 不可约，当且仅当 \(\mathcal{H}(\mathrm{G})\) 在 \(\mathcal{L}(\mathrm{E})\) 中的像关于逐点收敛拓扑是稠密的。（使用 IV, 第 326 页习题 24。）
\item
  令 H 为 G 的闭子群。令 \(\nu\) 为 G/H 上的拟不变测度（INT, VII, 第 56 页，第 2 节第 5 号定理 2），并令 \(\varrho\) 为与之关联的、取值于 \(\mathbf{R}_{+}^{*}\) 的 G 上函数。由下式定义的映射 \(\varrho(g)\)
\end{enumerate}

\[
\varrho (g) f (x) = \left(\frac {\varrho (g ^ {- 1} x)}{\varrho (x)}\right) ^ {1 / 2} f (g ^ {- 1} x)
\]

定义了 G 在 \(\mathcal{L}^{2}(\mathrm{G / H},\nu)\) 上的连续线性表示，并通过取商定义了 G 在 \(\mathrm{L}^2 (\mathrm{G / H},\nu)\) 上的酉表示，称为 \(\mathrm{G / H}\) 的拟正则表示。

\begin{enumerate}
\def\labelenumi{\arabic{enumi})}
\setcounter{enumi}{2}
\tightlist
\item
  令 G = R，令 H 为赋予离散拓扑并配备计数测度的群 R；该测度是哈尔测度。
\end{enumerate}

\begin{enumerate}
\def\labelenumi{\alph{enumi})}
\item
  对每个 \(t \in G\) 及每个 \(f \in L^{2}(H)\)，映射 \(\varrho(t)f: x \mapsto f(x + t)\) 属于 \(L^{2}(H)\)；映射 \(\varrho(t)\) 是 \(L^{2}(H)\) 上的酉映射。b) 对所有 \(f,g \in L^{2}(H)\)，由 \(t \mapsto \langle f | \varrho(t)g \rangle\) 定义的从 G 到 C 的映射是可测的。
\item
  令 \(f\) 和 \(g\) 属于 \(\mathrm{L}^2(\mathrm{H})\)。若存在 \(t \in \mathbf{R}\) 使得 \(\langle f | \varrho(t)g \rangle\) 非零，则映射 \(t \mapsto \langle f | \varrho(t)g \rangle\) 从 \(\mathbf{G}\) 到 \(\mathbf{C}\) 不连续。（特别地，E 可数生成这一假设不能从 INT, VIII, 第 171 页第 4 节第 7 号推论 2 中省略。）
\end{enumerate}

\(d)\) 存在 \(f \in \mathrm{L}^2(\mathrm{H})\)，使得 \(\mathrm{L}^2(\mathrm{H})\) 中由元素 \(\varrho(t)f\)（其中 \(t \in \mathbf{R}\)）生成的子空间在 \(\mathrm{L}^2(\mathrm{H})\) 中稠密。

\begin{enumerate}
\def\labelenumi{\arabic{enumi})}
\setcounter{enumi}{3}
\item
  令 \(p \geqslant 1\)。G 在 \(\mathrm{L}^{p}(\mathrm{G}, \mu)\) 上的双正则表示的核，是 G 的中心在同态 \(g \mapsto (g, g)\) 下的像。
\item
  令 G 为局部紧单模群。
\end{enumerate}

\begin{enumerate}
\def\labelenumi{\alph{enumi})}
\item
  若存在 G 的有限维平方可积表示，则 G 是紧的。
\item
  若存在 G 的平方可积表示，则 G 的中心是紧的。
\end{enumerate}

\begin{enumerate}
\def\labelenumi{\arabic{enumi})}
\setcounter{enumi}{5}
\tightlist
\item
  假设群 G 是单模的。记 \(\widehat{G}_{d} \subset \widehat{G}\) 为 G 的不可约酉平方可积表示的等价类集。对于 \(\pi \in \widehat{G}_{d}\)，记 \(p_{\pi}\) 为 \(L^{2}(G)\) 上的正交投影，其像是 G 的右正则表示的 \(\pi\)-同型分量。
\end{enumerate}

\(a)\) 令 \(f \in \mathcal{L}^2(\mathrm{G})\) 且 \(\pi \in \widehat{\mathrm{G}}_d\)。记 E 为 \(\pi\) 的表示空间。映射 \(g \mapsto f(g)\pi(g)\) 属于 \(\mathrm{L}^1(\mathrm{G};\mathcal{L}(\mathrm{E}))\)；自同态

\[
v = \int_ {\mathrm{G}} f (g) \pi (g) d \mu (g)
\]

是希尔伯特–施密特自同态，记为 \(v = \pi(f)\)。

\begin{enumerate}
\def\labelenumi{\alph{enumi})}
\setcounter{enumi}{1}
\item
  令 \(f \in \mathcal{L}^2(\mathrm{G})\)。有 \(\| \pi(\overline{f}) \|_2 = \| p_\pi(f) \|\)。
\item
  令 \(L_{d}^{2}(G)\) 为 \(L^{2}(G)\) 的闭子空间，它由 \(\pi\)-同型分量生成，这些分量属于 \(L^{2}(G)\)，且 \(\pi \in \widehat{G}_{d}\)。它是各正交投影 \(p_{\pi}\) 的像（对 \(\pi \in \widehat{G}_{d}\)）的希尔伯特直和。对于每个 \(f \in L_{d}^{2}(G)\)，有
\end{enumerate}

\[
\| f \| ^ {2} = \sum_ {\pi \in \widehat {\mathrm{G}} _ {d}} c (\pi) \| \pi (f) \| _ {2} ^ {2}
\]

其中 \(c(\pi)\) 是 \(\pi\) 相对于 \(\{e\}\) 的形式次数。

\begin{enumerate}
\def\labelenumi{\arabic{enumi})}
\setcounter{enumi}{6}
\tightlist
\item
  令 \(n \in N\)。确定 \(R^{n}\) 的酉自守表示（即对每个离散子群 \(\Gamma\)，它们都是 \(\Gamma\)-自守的 \(R^{n}\) 的表示）。
\end{enumerate}

8) 假设 G 是局部紧单模群。令 Γ 为 G 的离散子群，使商空间 X = Γ\textbackslash G 是紧的。令 μ 为 G 上的哈尔测度，β 为 Γ 上的计数测度；记 \(\widetilde{\mu} = \mu/\beta\)。

令 \(\varrho\) 为 V, 第 404 页例 3 中定义的 G 在 \(\mathrm{L}^{2}(\mathrm{X},\widetilde{\mu})\) 上的酉表示。

对于 \(\gamma\in\Gamma\)，记 \(\Gamma_{\gamma}\)（分别为 \(\mathbf{G}_{\gamma}\)）为 \(\gamma\) 在 \(\Gamma\)（分别在 \(\mathbf{G}\)）中的中心化子。\(a)\) 群 \(\mathbf{G}_{\gamma}\) 是单模的。

\begin{enumerate}
\def\labelenumi{\alph{enumi})}
\setcounter{enumi}{1}
\item
  令 \(\varphi\in\mathcal{K}(\mathrm{G})\)。对于每个 \(\Gamma\)-自守酉表示 \(\pi\)，自同态 \(\pi(\varphi)\) 都有有限迹。
\item
  有
\end{enumerate}

\[
\operatorname{Tr} (\varrho (\varphi)) = \sum_ {\pi \in \widehat {\mathrm{G}}} m _ {\pi} (\varrho) \operatorname{Tr} (\pi (\varphi))
\]

其中 \(m_{\pi}(\varrho)\) 是 \(\pi\) 在 \(\varrho\) 中的重数（“迹公式的谱侧”）。

\begin{enumerate}
\def\labelenumi{\alph{enumi})}
\setcounter{enumi}{3}
\tightlist
\item
  令 \(\Gamma^{\sharp}\) 为 \(\Gamma\) 的共轭类集。对于 \(\gamma \in \Gamma^{\sharp}\)，记 \(\beta_{\gamma}\) 为 \(\Gamma_{\gamma}\) 上的计数测度。有
\end{enumerate}

\[
\mathrm{Tr} (\varrho (\varphi)) = \sum_ {\gamma \in \Gamma^ {\sharp}} \int_ {\Gamma_ {\gamma} \backslash G} \varphi (x ^ {- 1} \gamma x) d (\mu / \beta_ {\gamma}) (x).
\]

\begin{enumerate}
\def\labelenumi{\alph{enumi})}
\setcounter{enumi}{4}
\tightlist
\item
  固定 \(\mu_{\gamma}\) 在 \(G_{\gamma}\) 上的哈尔测度，对所有 \(\gamma \in \Gamma\)。有
\end{enumerate}

\[
\operatorname{Tr} (\varrho (\varphi)) = \sum_ {\gamma \in \Gamma^ {\sharp}} \int_ {G _ {\gamma} \backslash G} \int_ {\Gamma_ {\gamma} \backslash G _ {\gamma}} \varphi (x ^ {- 1} y ^ {- 1} \gamma y x) d (\mu_ {\gamma} / \beta_ {\gamma}) (y) d (\mu / \mu_ {\gamma}) (x).
\]

\(f)\) 对每个 \(\gamma \in \Gamma\)，测度 \(\mu_{\gamma} / \beta_{\gamma}\) 在 \(\Gamma_{\gamma} \backslash G_{\gamma}\) 上是有限的。\(g)\) 有

\[
\mathrm{Tr} (\varrho (\varphi)) = \sum_ {\gamma \in \Gamma^ {\sharp}} (\mu_ {\gamma} / \beta_ {\gamma}) (\Gamma_ {\gamma} \backslash \mathrm{G} _ {\gamma}) \int_ {\mathrm{G} _ {\gamma} \backslash \mathrm{G}} \varphi (x ^ {- 1} \gamma x) d (\mu / \mu_ {\gamma}) (x)
\]

（“迹公式的几何侧”；将此表达式与 c) 中所得表达式相等联系起来的公式，称为 Γ-自守表示的“迹公式”。）

\begin{enumerate}
\def\labelenumi{\alph{enumi})}
\setcounter{enumi}{7}
\tightlist
\item
  置 G = R 且 \(\Gamma = Z\)。验证此时迹公式等价于泊松公式（参见 II, 第 239 页）。
\end{enumerate}

\begin{enumerate}
\def\labelenumi{\arabic{enumi})}
\setcounter{enumi}{8}
\item
  假设 G 是无紧因子的完美实半单李群。G 没有非平凡的有限维酉表示。
\item
  \begin{enumerate}
  \def\labelenumii{\alph{enumii})}
  \tightlist
  \item
    G 在 \(L^{2}(G)\) 上的左正则表示的每个矩阵系数都属于 \(\mathcal{C}_{0}(G)\)。
  \end{enumerate}
\end{enumerate}

\begin{enumerate}
\def\labelenumi{\alph{enumi})}
\setcounter{enumi}{1}
\tightlist
\item
  若 G 非紧，则 G 在 L²(G) 上的左正则表示不含非零有限维子表示。
\end{enumerate}

\begin{enumerate}
\def\labelenumi{\arabic{enumi})}
\setcounter{enumi}{10}
\tightlist
\item
  令 E 为希尔伯特空间。
\end{enumerate}

\begin{enumerate}
\def\labelenumi{\alph{enumi})}
\item
  令 \(\varrho\) 为群 \(\mathbf{R}\) 在 E 中的酉表示。令 D 为 E 的稠密子空间，在 \(\varrho\) 下稳定，并且对每个 \(x\in\mathrm{D}\)，映射 \(\psi_{x}:t\mapsto\varrho(t)x\) 在 R 上可微。令 u 为定义域为 D 且满足 \(u(x) = i^{-1}\psi_x'(0)\)（对所有 \(x \in D\)）的部分算子。则 u 本质自伴，且 \(\overline{u}\) 是 \(\varrho\) 的无穷小生成元。
\item
  令 u 为 E 上的自伴部分算子，D 为 \(\text{dom}(u)\) 的稠密子空间。假设对每个 \(t \in R\)，有 \(\exp(itu)(D) \subset D\)。则空间 D 在 \(E_{u}\) 中稠密。
\item
  令 u 为 E 上的自伴部分算子。令 F 为希尔伯特空间，v 为 F 上的自伴部分算子。部分算子 \(u \otimes 1_{F} + 1_{E} \otimes E\) 在 \(E \otimes F\) 上定义了 \(E \widehat{\otimes}_{2} F\) 上的一个本质自伴部分算子。
\end{enumerate}

\begin{enumerate}
\def\labelenumi{\arabic{enumi})}
\setcounter{enumi}{11}
\tightlist
\item
  令 G 为实李群，令 \(\varrho\) 为 G 在希尔伯特空间 E 中的酉表示。若元素 \(x \in \mathrm{E}\) 为 \(\mathrm{C}^{\infty}\) 类，亦即从 G 到 E 的映射 \(g \mapsto \varrho(g)x\) 为 \(\mathrm{C}^{\infty}\) 类，则称其为类向量。a) 令 \(\pi\) 为 G 在 F 中的酉表示，且 \(u: \mathrm{E} \to \mathrm{F}\) 为 G-态射。对于每个 \(x \in \mathrm{E}\) 的 \(\mathrm{C}^{\infty}\) 类向量，向量 \(u(x)\) 在 F 中也是 \(\mathrm{C}^{\infty}\) 类的。
\end{enumerate}

记 \(\mathrm{C}^{\infty}(\varrho)\) 为 E 中 \(C^{\infty}\) 类向量的空间。由问题 a)，每个从 \(\varrho\) 到 \(\pi\) 的 G-态射通过取子空间定义了从 \(\mathrm{C}^{\infty}(\varrho)\) 到 \(\mathrm{C}^{\infty}(\pi)\) 的线性映射。

\begin{enumerate}
\def\labelenumi{\alph{enumi})}
\setcounter{enumi}{1}
\tightlist
\item
  空间 \(\mathrm{C}^{\infty}(\varrho)\) 在 E 中稠密。
\end{enumerate}

\begin{enumerate}
\def\labelenumi{\arabic{enumi})}
\setcounter{enumi}{12}
\tightlist
\item
  令 \(\pi\) 为 G 在复希尔伯特空间 E 中的不可约酉表示。令 \(\varrho\) 为 G 在复希尔伯特空间 F 中的酉表示。假设存在一个从 E 到 F 的闭稠定部分算子 \(u\)，其定义域是 E 的 G-稳定子空间，并且满足 \(u \circ \pi(g) = \varrho(g) \circ u\) 对所有 \(g \in \mathbf{G}\) 成立。
\end{enumerate}

部分算子 \(u\) 属于 \(\mathcal{L}(\mathrm{E};\mathrm{F})\)，并且存在 \(\lambda > 0\)，使得 \(\lambda u\) 是等距的。

\begin{enumerate}
\def\labelenumi{\arabic{enumi})}
\setcounter{enumi}{13}
\tightlist
\item
  令 Z 为 G 的中心 C 的子群，且 C/Z 是紧的；令 \(\nu\) 为 G/Z 上的哈尔测度。令 \(\pi\) 为 G 在复希尔伯特空间 E 中的不可约酉表示。假设 \(\pi\) 模中心平方可积。令 \(u \in \mathcal{L}(\mathrm{E})\) 为核型自同态（IV, 第 169 页定义 4）。
\end{enumerate}

\begin{enumerate}
\def\labelenumi{\alph{enumi})}
\tightlist
\item
  对 E 中的每个 \(x\) 和 \(y\)，函数
\end{enumerate}

\[
g \mapsto \langle x | \pi (g) u \pi (g ^ {- 1}) y \rangle
\]

属于 \(\mathcal{F}(\mathrm{G / Z})\)。

\begin{enumerate}
\def\labelenumi{\alph{enumi})}
\setcounter{enumi}{1}
\tightlist
\item
  对 E 中的每个 \(x\) 和 \(y\)，有
\end{enumerate}

\[
\int_ {\mathrm{G} / \mathrm{Z}} \langle x \mid \pi (g) u \pi (g ^ {- 1}) y \rangle d \nu (g) = \frac {1}{c _ {\mathrm{Z}} (\pi)} \langle x \mid y \rangle \operatorname{Tr} (u).
\]

\begin{enumerate}
\def\labelenumi{\arabic{enumi})}
\setcounter{enumi}{14}
\tightlist
\item
  令 Z 为 G 的中心 C 的子群，且 C/Z 是紧的；令 \(\nu\) 为 G/Z 上的哈尔测度。令 \(\pi\) 为 G 在复希尔伯特空间 E 中的不可约酉表示。记 \(\chi\) 为 \(\pi\) 的中心特征标在 Z 上的限制。
\end{enumerate}

记 \(f_{x,y}\) 为矩阵系数 \(g \mapsto \langle x | \pi(g)y \rangle\)，其中 \((x, y) \in \mathrm{E}^{2}\)。令 \(p \geqslant 1\) 为实数，令 \(q\) 为 p 的共轭指数。

\begin{enumerate}
\def\labelenumi{\alph{enumi})}
\item
  令 \(y \in \mathrm{E}\)，使得矩阵系数 \(f_{x,y}\) 属于 \(\mathcal{L}_{\chi}^{p}(\mathrm{G})\)。映射 \(x \mapsto f_{x,y}\) 是从 \(\pi\) 到 \(\boldsymbol{\delta}_{\mathrm{G},\chi}^{(p)}\) 的连续 G-态射。
\item
  令 \(\varrho\) 为 G 在复希尔伯特空间 F 中的不可约酉表示，其中心特征标限制所得的 \(\chi\) 定义在 \(Z\) 上。假设 \(p > 1\)，且 \(\pi\) 的矩阵系数属于 \(\mathcal{L}_{\chi}^{p}(\mathrm{G})\)，\(\varrho\) 的矩阵系数属于 \(\mathcal{L}_{\chi}^{q}(\mathrm{G})\)。令 \(x_0 \in \mathrm{E}\) 且 \(y_0 \in \mathrm{F}\)。对于每个 \(x \in \mathrm{E}\)，存在从 F 到 C 的连续线性泛函 \(\ell_x\)，使得
\end{enumerate}

\[
\ell_ {x} (y) = \int_ {\mathrm{G/Z}} \overline {{\langle x _ {0} | \pi (g) x \rangle}} \langle y _ {0} | \varrho (g) y \rangle d \nu (g)
\]

对所有 \(y \in \mathbf{F}\) 都成立。

\begin{enumerate}
\def\labelenumi{\alph{enumi})}
\setcounter{enumi}{2}
\item
  存在从 E 到 F 的连续线性映射 u，使得有 \(\ell_{x}(y)=\langle u(x)\mid y\rangle\)，对所有 \((x,y)\in\mathrm{E}\times\mathrm{F}\) 成立。
\item
  有
\end{enumerate}

\[
\int_ {\mathrm{G} / \mathrm{Z}} \overline {{\langle x | \pi (g) x ^ {\prime} \rangle}} \langle y | \varrho (g) y ^ {\prime} \rangle d \nu (g) = 0
\]

对 \((x,x^{\prime},y,y^{\prime})\in \mathrm{E}^{2}\times \mathrm{F}^{2}\) 都成立。

\begin{enumerate}
\def\labelenumi{\alph{enumi})}
\setcounter{enumi}{4}
\tightlist
\item
  假设 p = 1。\(\pi\) 的矩阵系数与 \(\varrho\) 的矩阵系数正交。
\end{enumerate}

\begin{enumerate}
\def\labelenumi{\arabic{enumi})}
\setcounter{enumi}{15}
\tightlist
\item
  令 \(\mathbf{G} = \mathbf{S}\mathbf{L}(2, \mathbf{R})\)，令 H 为 C 中由满足 \(\mathcal{I}(z) > 0\) 的复数 z 组成的开子集。令 \(\mu\) 为 C 上的勒贝格测度。对于 G 中的每个元素
\end{enumerate}

\[
g = \left( \begin{array}{c c} a & b \\ c & d \end{array} \right)
\]

及每个 \(z \in \mathbf{H}\)，置 \(j(g, z) = cz + d\)。

\begin{enumerate}
\def\labelenumi{\alph{enumi})}
\tightlist
\item
  群 G 通过下式在 H 上连续地左作用
\end{enumerate}

\[
\left( \begin{array}{c c} a & b \\ c & d \end{array} \right) \cdot z = \frac {a z + b}{c z + d}.
\]

测度 \(\nu = y^{-2} \cdot \mu\) 是 G-不变的。

\begin{enumerate}
\def\labelenumi{\alph{enumi})}
\setcounter{enumi}{1}
\tightlist
\item
  令 n 为满足 \(n \geqslant 2\) 的整数。记 \(E_{n}\) 为全纯函数 \(f: H \to C\) 所成的向量空间，其中函数 \(z \mapsto |f(z)|^{2} \mathcal{I}(z)^{n}\) 属于 \(\mathcal{L}^{2}(\mathbf{H}, \nu)\)。空间 \(E_{n}\) 在赋予半双线性形式
\end{enumerate}

\[
\langle f \mid g \rangle = \int_ {\mathbf {H}} \overline {{f (z)}} g (z) \mathcal {I} (z) ^ {n} d \nu (z)
\]

后，是可数型且无限维的希尔伯特空间。

\begin{enumerate}
\def\labelenumi{\alph{enumi})}
\setcounter{enumi}{2}
\tightlist
\item
  对每个 \(g \in G\) 及每个 \(f \in E_{n}\)，置
\end{enumerate}

\[
\pi_ {n} (g) f (z) = j (g ^ {- 1}, z) ^ {- n} f (g ^ {- 1} \cdot z).
\]

映射 \(\pi_n\) 定义了 G 在 \(\mathrm{E}_n\) 中的酉表示。\(d)\) 由函数 \(f_n\) 定义为 \(f_n(z) = (z + i)^{-n}\)，它属于 \(\mathrm{E}_n\)；它满足 \(\pi_n(k)f_n = \chi_{-n}(k)f_n\) 对所有 \(k \in \mathbf{SO}(\mathbf{R}^2)\) 成立，其中 \(\chi_n\) 是满足下式的特征标

\[
\chi_ {n} \Big (\left( \begin{array}{c c} \cos (\theta) & \sin (\theta) \\ - \sin (\theta) & \cos (\theta) \end{array} \right) \Big) = e ^ {i n \theta}
\]

对所有 \(\theta\in\mathbf{R}\) 都成立。

\begin{enumerate}
\def\labelenumi{\alph{enumi})}
\setcounter{enumi}{4}
\item
  令 F 为 \(E_{n}\) 的子表示。\(E_{n}\) 的正交投影，其像是 \(\chi_{-n}\)-同型分量，属于将 \(\pi_{n}\) 限制到 \(\mathbf{SO}(\mathbf{R}^{2})\) 后的表示，由 \(f \mapsto -(2i)^{n}f(i)f_{n}\) 给出。（使用 V, 第 465 页推论 2 以及柯西积分公式，FVC，即将出版。）
\item
  表示 \(\pi_n\) 是不可约的。
\end{enumerate}

\(g)\) 矩阵系数 \(g \mapsto \langle f_n | \pi_n(g)f_n \rangle\) 非零且属于 \(\mathcal{L}^2(\mathrm{G})\)。

\begin{enumerate}
\def\labelenumi{\alph{enumi})}
\setcounter{enumi}{7}
\tightlist
\item
  表示 \(\pi_{n}\) 是平方可积的，并且存在常数 c > 0，使得 \(\pi_{n}\) 的形式次数为 \(c(n - 1)\)。
\end{enumerate}

\begin{enumerate}
\def\labelenumi{\arabic{enumi})}
\setcounter{enumi}{16}
\item
  令 G 为连通实李群。G 存在忠实有限维酉表示，当且仅当群 G 局部同构于紧群。
\item
  令 X 为赋有有界正测度 \(\mu\) 的局部紧空间。对于每个巴拿赫空间 F，记 S(X, \(\mu\) ; F) 为从 X 到 F 的 \(\mu\)-可测函数按 \(\mu\)-几乎处处相等关系所得的等价类空间，并赋予依测度收敛拓扑（INT, IV, 第 194 页，第 5 节第 11 号）。记 A 为 \(\mathcal{L}(\mathrm{L}_{\mathbf{C}}^{2}(\mathrm{X},\mu))\) 的交换子代数，它由函数 \(f \in \mathcal{L}_{\mathbf{C}}^{\infty}(\mathrm{X},\mu)\) 所对应的乘法自同态组成。
\end{enumerate}

记 \(e\colon \mathbf{R}\to \mathbf{C}\) 为同态 \(t\mapsto \exp (2i\pi t)\)。

令 E 为拓扑向量空间。

\begin{enumerate}
\def\labelenumi{\alph{enumi})}
\tightlist
\item
  令 \(f \in \mathrm{S}(\mathrm{X}, \mu; \mathbf{R})\)，使得 \(|f(x)| \geqslant 1\) 对 \(\mu\)-几乎所有 \(x \in \mathrm{X}\) 都成立。存在实数 \(t \in ]0, 1[\)，使得
\end{enumerate}

\[
\mu (\{x \in \mathrm{X} \mid t f (x) \in [ 1 / 4, 3 / 4 ] + \mathbf {Z} \}) \geqslant \frac {2}{5} \mu (\mathrm{X}).
\]

\begin{enumerate}
\def\labelenumi{\alph{enumi})}
\setcounter{enumi}{1}
\item
  令 u 为从 E 到 \(\mathrm{S}(\mathrm{X},\mu;\mathbf{R})\) 的线性映射。对于每个 \(x\in E\)，记 \(\varrho_{u}(x)\) 为 A 中由函数 \(e\circ u(x)\) 的乘法自同态给出的元素，该函数属于 \(\mathcal{L}_{\mathbf{C}}^{\infty}(\mathrm{X},\mu)\)。映射 \(\varrho_{u}\) 是从 G 到 \(\mathrm{L}_{\mathbf{C}}^{2}(\mathrm{X},\mu)\) 的连续表示，当且仅当映射 u 连续。（若 u 不连续，使用上一问题验证 \(\varrho\) 不连续。）
\item
  令 \(\varrho\) 为 E 在 \(\mathrm{L}^2 (\mathrm{X},\mu)\) 中的酉表示，且其像包含于 \(\mathcal{A}\)。存在唯一的连续线性映射 \(u\colon \mathrm{E}\to \mathrm{S}(\mathrm{X},\mu ;\mathbf{R})\)，使得 \(\varrho = \varrho_{u}\)。（先考虑 \(\mathrm{E} = \mathbf{R}\) 的情形，然后应用斯通定理。）
\end{enumerate}

\begin{enumerate}
\def\labelenumi{\arabic{enumi})}
\setcounter{enumi}{18}
\item
  令 G 为有限群，H 为 G 的子群。对于 H 在有限维复希尔伯特空间 E 中的每个酉表示 \(\pi\)，证明诱导酉表示 \(\operatorname{Ind}_{\mathrm{H}}^{\mathrm{G}}(\pi)\) 同构于 A, VIII, 第 392 页第 21 节第 3 号中定义的诱导表示。
\item
  令 G 为无穷远处可数的局部紧群，令 H 为 G 的闭子群。令 \(\pi\) 为 H 在希尔伯特空间 E 中的酉表示，令 \(\varrho\) 为由 \(\pi\) 诱导的 G 的酉表示。记 F 为 \(\varrho\) 的表示空间。
\end{enumerate}

在 G 上固定哈尔测度 \(\mu\)，在 H 上固定哈尔测度 \(\beta\)，并记 \(\kappa\colon\mathbf{G}\to\mathbf{R}_{+}^{*}\) 为满足 \(\kappa(xh)=\Delta_{\mathrm{H}}(h)\Delta_{\mathrm{G}}(h)^{-1}\kappa(x)\) 对 \((x,h)\in\mathbf{G}\times\mathrm{H}\) 成立，且满足 \(\kappa(e)=1\) 的连续函数。

\begin{enumerate}
\def\labelenumi{\alph{enumi})}
\item
  若 G/H 具有有界 G-不变测度，且存在非零的 H-不变向量 \(x \in E\)，则存在非零的 G-不变向量 \(y \in F\)。
\item
  反之，设 \(y \in F\) 是 \(\varrho\) 的非零 G-不变向量。证明存在非零函数 \(f \in L_{\overline{\pi}}^{2}\)，使得
\end{enumerate}

\[
f (x) = \left(\frac {\kappa (g ^ {- 1} x)}{\kappa (x)}\right) ^ {1 / 2} f (g ^ {- 1} x)
\]

对所有 \((g, x) \in \mathrm{G}^{2}\) 都成立。推出 \(\kappa = 1\)，从而 G/H 上的测度 \(\mu/\beta\) 有定义且有界，最后证明 \(f(e) \in \mathrm{E}\) 是非零 H-不变向量。
% semantic-fragments: legacy-input-seam
\relax{}
